\documentclass[11pt]{letter}
\usepackage[margin=1in]{geometry}
\usepackage{hyperref}
\usepackage{xcolor}

% Author information
\signature{Jiahao Wu \\ Corresponding Author}
\address{The University of Hong Kong\\
Hong Kong, China\\
Email: 13620926353@163.com}

\date{February 27, 2026}

% Make closing and signature left-aligned
\longindentation=0pt

\begin{document}

\begin{letter}{%
The Editor\\
\textit{Robotica}\\
Cambridge University Press}

\opening{Dear Editor,}

We are pleased to submit our revised manuscript \textbf{ROB-2026-0021}, entitled:

\begin{quote}
\textbf{``Cross-Platform Learnable Fuzzy Gain-Scheduled Proportional-Integral-Derivative Tuning via Physics-Constrained Meta-Learning and Reinforcement Learning Adaptation''}
\end{quote}

\noindent for further consideration for publication as a Research Article in \textit{Robotica}.

We thank the Editor and the anonymous reviewers for their thorough and constructive evaluations. Their feedback has substantially improved the rigor and clarity of this work. We have carefully addressed \textbf{all nine reviewer comments} in the accompanying point-by-point response letter.

\medskip
\noindent\textbf{Summary of Key Revisions:}

\begin{enumerate}
    \item \textbf{Sensitivity \& justification analyses (Reviewer 1, Comments 1--3):} Added Table~1 (controllability threshold sensitivity) and Table~2 (optimizer pilot comparison with CMA-ES, L-BFGS-B, PSO baselines), together with extended discussion of the data filtering rationale and hybrid DE+NM optimization choice.

    \item \textbf{Generalization evidence (Reviewer 1, Comment 4; Reviewer 2, Comment 3):} Added leave-one-platform-out analysis (Table~7) showing that Laikago---contributing $\leq$0.4\% of training data---achieves 5.91° MAE via feature-space generalization alone, constituting a \textit{de facto} cross-morphology transfer test.

    \item \textbf{Optimization ceiling effect (Reviewer 2, Comments 1--2):} Introduced formal per-joint RL gain rate metric $G_{\mathrm{RL},j}$, added Table~10 (quantitative characterization) and new Figure~7 (scatter plot with Spearman correlation $\rho_s = 0.38$), providing rigorous characterization of when RL adaptation is beneficial versus unnecessary.

    \item \textbf{Computational budget transparency (Reviewer 2, Comment 4):} Added Appendix~C with unified resource breakdown table distinguishing one-time offline costs from per-platform deployment costs; clarified amortization of meta-learning overhead.

    \item \textbf{Manual Fuzzy-PID baseline (Reviewer 1, Comment 5):} Added expert-tuned Fuzzy-PID baseline to Table~3, providing direct comparison with interpretable scheduling approaches.

    \item \textbf{RL adaptation subset ablation (Reviewer 2, Comment 5):} Added Appendix~D with three-configuration ablation study (scales-only vs.\ consequents vs.\ combined), demonstrating that the proposed scales-only adaptation achieves the best balance of performance (+16.6\%) and stability (0/5 divergence).
\end{enumerate}

\noindent All revisions are highlighted in {\color{blue}blue} in the manuscript for easy identification.

\medskip
\noindent\textbf{Novelty and Key Findings:}

This work presents a hierarchical framework that eliminates the need for per-platform expert tuning of PID controllers in robotics. The key novelty lies in combining physics-constrained virtual robot synthesis (generating 232 valid training variants from only 3 base platforms) with meta-learning for instant zero-shot initialization and lightweight RL for online adaptation. Validated on morphologically diverse platforms (9-DOF manipulator and 12-DOF quadruped), the method achieves up to 80.4\% error reduction on challenging joints and 19.2\% improvement under parameter uncertainty, while reducing deployment time from hours of expert tuning to under 10 minutes of automated adaptation. The newly identified \textit{optimization ceiling effect} provides practitioners with a principled criterion for deciding when RL refinement is cost-effective.

\medskip
\noindent\textbf{Declarations:} This is original work not under consideration elsewhere. No competing interests exist and no external funding was received. Complete author contributions, data availability, and AI-assisted technology declarations are provided in the manuscript. All code and data will be publicly available upon acceptance under the MIT License.

\medskip
We hope the revised manuscript now meets the standards of \textit{Robotica}. We look forward to your editorial decision.

\closing{Sincerely,}

\end{letter}

\end{document}
